\documentclass[11pt]{article}
\usepackage[T1]{fontenc}
\usepackage{lmodern}
\usepackage{amsmath,amssymb,amsthm,mathtools}
\usepackage{booktabs}
\usepackage[margin=1in]{geometry}
\usepackage[hidelinks]{hyperref}
\usepackage{microtype}

\newtheorem{theorem}{Theorem}
\newcommand{\D}{\mathbb D}
\newcommand{\A}{\mathcal A}
\newcommand{\M}{\mathcal M}

\title{The multiplier algebra of the coefficient-weighted Bergman scale\\
\large A literature-implied full answer to the question in arXiv:2105.07760}
\author{Autonomous literature and proof audit, lane 18}
\date{11 August 2026}

\begin{document}
\maketitle

\begin{center}
\fbox{\parbox{0.91\textwidth}{\textbf{Status.}
Full literature-implied answer at the level of an exact function-theoretic
criterion. Nothing here is claimed as new: the ingredients predate the source
paper.}}
\end{center}

\section{The source question}

Gallardo-Guti\'errez and Partington~\cite{source} define, for
$\alpha\in\mathbb R$,
\[
 \A_\alpha=\left\{f(z)=\sum_{n\geq0}a_nz^n:
 \|f\|_\alpha^2:=\sum_{n\geq0}|a_n|^2(n+1)^\alpha<\infty\right\},
\]
and let $\M_\alpha$ be the analytic functions $g$ for which
$f\mapsto gf$ is bounded on $\A_\alpha$. On source PDF page 3 they write:
\begin{quote}
For other instances of $\alpha$ this is no longer the case, and it is an open
question to identify them for any $\alpha$.
\end{quote}
Here ``the case'' refers to the equality with $H^\infty$ at $\alpha=-1,0$.
The literature already supplies a complete phase diagram, although the middle
range is naturally described by a Carleson condition rather than by a simpler
named algebra.

\section{Complete phase diagram}

A positive measure $\mu$ on $\D$ is called a Carleson measure for $\A_\alpha$
if
\[
 \int_\D |f|^2\,d\mu\leq C\|f\|_\alpha^2
 \qquad(f\in\A_\alpha).
\]

\begin{theorem}\label{thm:phase}
For every real $\alpha$, the multiplier algebra of $\A_\alpha$ is as follows:
\[
\begin{array}{ccl}
\alpha\leq0
 &:& \M_\alpha=H^\infty,\\[2mm]
0<\alpha\leq1
 &:& \displaystyle
 \M_\alpha=\left\{g\in H^\infty:
 |g'(z)|^2(1-|z|^2)^{1-\alpha}\,dA(z)
 \text{ is Carleson for }\A_\alpha\right\},\\[3mm]
\alpha>1
 &:& \M_\alpha=\A_\alpha.
\end{array}
\]
The identities are set identities, and the corresponding natural norms are
equivalent.
\end{theorem}

The first and third regimes occur already in Taylor's one-variable multiplier
theory~\cite{Taylor}. Jupiter and Redett~\cite{JR}, whose paper explicitly
extends Taylor's results, record the stronger polydisc versions as Proposition
3.9 and Theorem 3.10 on physical PDF page 13. The middle regime is the classical
Kerman--Sawyer criterion~\cite{KS}, quoted explicitly by Bao, Lou, Qian, and
Wulan~\cite[p.~1722; physical PDF p.~23]{Bao}. The classical Dirichlet endpoint
$\alpha=1$ is also covered by Stegenga's multiplier theorem~\cite{Stegenga}.

\section{Parameter translation and direct verification}

For $0<\alpha\leq1$, beta-function asymptotics give the norm equivalence
\begin{equation}\label{eq:derivative}
 \|f\|_\alpha^2\asymp |f(0)|^2+
 \int_\D |f'(z)|^2(1-|z|^2)^{1-\alpha}\,dA(z).
\end{equation}
Thus the source space is the standard weighted Dirichlet space $\mathcal D_K$
with $K(t)=t^{1-\alpha}$, up to an equivalent norm. Since
$0<\alpha\leq1$, this $K$ is increasing and concave and
$t/K(t)=t^\alpha\to0$. The Kerman--Sawyer theorem therefore applies. Replacing
$1-|z|$ by $1-|z|^2$ changes the weight only by fixed constants.

For completeness, the multiplier criterion itself follows immediately from
\eqref{eq:derivative}. If $M_g$ is bounded, then the reproducing-kernel identity
$M_g^*k_z=\overline{g(z)}k_z$ gives $g\in H^\infty$. Moreover,
\[
 fg'=(fg)'-gf'
\]
shows that $|g'|^2(1-|z|^2)^{1-\alpha}dA$ is Carleson. Conversely, if $g$ is
bounded and that measure is Carleson, then
\[
 (fg)'=f g'+g f'
\]
and \eqref{eq:derivative} give $\|fg\|_\alpha\lesssim\|f\|_\alpha$.

The endpoint regimes are equally direct. For $\alpha<0$, $\A_\alpha$ has an
equivalent radial weighted-Bergman norm, so $H^\infty\subseteq\M_\alpha$; the
kernel argument gives the reverse inclusion. The case $\alpha=0$ is $H^2$.
If $\alpha>1$, weighted Cauchy--Schwarz gives $\ell^1$ summability of the Taylor
coefficients. Together with
\[
 (m+n+1)^{\alpha/2}\lesssim
 (m+1)^{\alpha/2}+(n+1)^{\alpha/2},
\]
Young's convolution inequality proves that $\A_\alpha$ is a Banach algebra.
Hence $\A_\alpha\subseteq\M_\alpha$, while the reverse inclusion follows by
applying a multiplier to $1$.

\section{Provenance verdict}

The three regimes exhaust $\mathbb R$, so the source sentence has a full answer
at the criterion level. The status is nevertheless \emph{literature-implied},
not ``later paper explicitly answered'': the cited theorems predate
arXiv:2105.07760 and require the norm/parameter identification above. This
packet should not be counted as an original solve.

\small
\begin{thebibliography}{9}
\bibitem{source}
E. A. Gallardo-Guti\'errez and J. R. Partington,
\emph{Multiplication by a finite Blaschke product on weighted Bergman spaces:
commutant and reducing subspaces}, arXiv:2105.07760, especially PDF p.~3.

\bibitem{Taylor}
G. D. Taylor, \emph{Multipliers on $D_\alpha$},
Trans. Amer. Math. Soc. \textbf{123} (1966), 229--240,
DOI: 10.1090/S0002-9947-1966-0206696-6.

\bibitem{Stegenga}
D. A. Stegenga, \emph{Multipliers of the Dirichlet space},
Illinois J. Math. \textbf{24} (1980), 113--139.

\bibitem{KS}
R. Kerman and E. Sawyer,
\emph{Carleson measures and multipliers of Dirichlet-type spaces},
Trans. Amer. Math. Soc. \textbf{309} (1988), 87--98,
DOI: 10.1090/S0002-9947-1988-0957062-1.

\bibitem{JR}
D. Jupiter and D. Redett, \emph{Multipliers on Dirichlet type spaces},
Acta Sci. Math. (Szeged) \textbf{72} (2006), 179--203;
arXiv:math/0510193.

\bibitem{Bao}
G. Bao, Z. Lou, R. Qian, and H. Wulan,
\emph{On multipliers of Dirichlet type spaces},
Complex Anal. Oper. Theory \textbf{9} (2015), 1701--1732,
DOI: 10.1007/s11785-015-0444-0.
\end{thebibliography}

\end{document}
